\documentclass[a4paper,11pt]{article}

\usepackage[utf8]{inputenc}
\usepackage[T1]{fontenc}
\usepackage{geometry}
\usepackage{xcolor}
\usepackage{listings}
\usepackage[hidelinks]{hyperref}
\usepackage{enumitem}

% Page Geometry setup
\geometry{margin=1in}

% Color definitions for code listings
\definecolor{codegreen}{rgb}{0,0.6,0}
\definecolor{codegray}{rgb}{0.5,0.5,0.5}
\definecolor{codepurple}{rgb}{0.58,0,0.82}
\definecolor{backcolour}{rgb}{0.95,0.95,0.92}

% JSON language definition for listings (so [language=json] works)
\lstdefinelanguage{json}{
  basicstyle=\ttfamily\footnotesize,
  numbers=left,
  numberstyle=\tiny\color{codegray},
  stepnumber=1,
  numbersep=5pt,
  showstringspaces=false,
  breaklines=true,
  frame=single,
  backgroundcolor=\color{backcolour},
  literate=
   *{0}{{{\color{codepurple}0}}}{1}
    {1}{{{\color{codepurple}1}}}{1}
    {2}{{{\color{codepurple}2}}}{1}
    {3}{{{\color{codepurple}3}}}{1}
    {4}{{{\color{codepurple}4}}}{1}
    {5}{{{\color{codepurple}5}}}{1}
    {6}{{{\color{codepurple}6}}}{1}
    {7}{{{\color{codepurple}7}}}{1}
    {8}{{{\color{codepurple}8}}}{1}
    {9}{{{\color{codepurple}9}}}{1}
    {:}{{{\color{magenta}{:}}}}{1}
    {,}{{{\color{magenta}{,}}}}{1}
    {\{}{{{\color{magenta}{\{}}}}{1}
    {\}}{{{\color{magenta}{\}}}}}{1}
    {[}{{{\color{magenta}{[}}}}{1}
    {]}{{{\color{magenta}{]}}}}{1}
}

% Listings configuration (general)
\lstdefinestyle{mystyle}{
    backgroundcolor=\color{backcolour},
    commentstyle=\color{codegreen},
    keywordstyle=\color{magenta},
    numberstyle=\tiny\color{codegray},
    stringstyle=\color{codepurple},
    basicstyle=\ttfamily\footnotesize,
    breakatwhitespace=false,
    breaklines=true,
    captionpos=b,
    keepspaces=true,
    numbers=left,
    numbersep=5pt,
    showspaces=false,
    showstringspaces=false,
    showtabs=false,
    tabsize=2,
    frame=single
}
\lstset{style=mystyle}

\title{\textbf{Controller Documentation}}

\begin{document}
\maketitle

\section*{Overview}
This document describes the endpoints of seven \texttt{controller.java} files, including:

\begin{enumerate}[label=\arabic*.]
    \item AnalyticsController endpoints
    \item AuthController endpoints
    \item BundleController endpoints
    \item CategoryController endpoints
    \item GamificationController endpoints
    \item IssueController endpoints
    \item ReservationController endpoints
\end{enumerate}

\hrulefill

\section{AnalyticsController}
\textbf{Base path:} \texttt{/api/analytics}

\subsection{Get seller dashboard metrics}
\texttt{GET /api/analytics/dashboard/\{sellerId\}}

\paragraph{Example Response:}
\begin{lstlisting}[language=json]
{
  "sellerName": "ABC",
  "totalBundlesPosted": 42,
  "totalQuantity": 180,
  "collectedCount": 120,
  "noShowCount": 15,
  "expiredCount": 5,
  "sellThroughRate": 66.6,
  "wasteAvoidedGrams": 36000,
  "openIssueCount": 2
}
\end{lstlisting}

\subsection{Get seller sell-through breakdown}
\texttt{GET /api/analytics/sell-through/\{sellerId\}}

\paragraph{Example Response:}
\begin{lstlisting}[language=json]
{
  "collected": 120,
  "noShow": 15,
  "expired": 5,
  "collectionRate": 85.7,
  "noShowRate": 10.7
}
\end{lstlisting}

\paragraph{Notes:}
\begin{itemize}
    \item \texttt{collectionRate = collected / total * 100}
    \item \texttt{noShowRate = noShow / total * 100}
\end{itemize}

\subsection{Get seller waste metrics (with assumptions)}
\texttt{GET /api/analytics/waste/\{sellerId\}}

\paragraph{Example Response:}
\begin{lstlisting}[language=json]
{
  "bundlesCollected": 120,
  "weightSavedGrams": 36000,
  "co2eSavedKg": 90.0,
  "mealsSaved": 120,
  "weightAssumption": "Average bundle weight: 300g",
  "co2eAssumption": "CO2e factor: 2.5 kg CO2e per kg food waste avoided",
  "mealsAssumption": "Meals estimate: 1 meal per bundle"
}
\end{lstlisting}

\hrulefill

\section{AuthController}
\textbf{Base path:} \texttt{/api/auth}

\subsection{Register a new user (and optionally create a Seller record)}
\texttt{POST /api/auth/register}

\paragraph{Example Request:}
\begin{lstlisting}[language=json]
{
  "email": "seller@example.com",
  "password": "PlainTextPasswordHere",
  "role": "SELLER",
  "businessName": "ABC",
  "location": "Exeter"
}
\end{lstlisting}

\paragraph{Example Response:}
\begin{lstlisting}[language=json]
{
  "token": "eyJhbGciOiJIUzI1NiIsInR5cCI6IkpXVCJ9....",
  "userId": "2f3a2c7d-6a1b-4e6a-9c8b-1f2d3c4b5a6e",
  "email": "seller@example.com",
  "role": "SELLER"
}
\end{lstlisting}

\paragraph{Notes:}
\begin{itemize}
    \item If role is \texttt{SELLER} and \texttt{businessName != null}, it creates a \texttt{Seller} linked to this user.
    \item Returns a JWT token in response.
\end{itemize}

\subsection{Login (get JWT)}
\texttt{POST /api/auth/login}

\paragraph{Example Request:}
\begin{lstlisting}[language=json]
{
  "email": "seller@example.com",
  "password": "PlainTextPasswordHere"
}
\end{lstlisting}

\paragraph{Example Response:}
\begin{lstlisting}[language=json]
{
  "token": "eyJhbGciOiJIUzI1NiIsInR5cCI6IkpXVCJ9....",
  "userId": "2f3a2c7d-6a1b-4e6a-9c8b-1f2d3c4b5a6e",
  "email": "seller@example.com",
  "role": "SELLER"
}
\end{lstlisting}

\subsection{Get current user profile (from SecurityContext)}
\texttt{GET /api/auth/me}

\paragraph{Example Response:}
\begin{lstlisting}[language=json]
{
  "token": null,
  "userId": "2f3a2c7d-6a1b-4e6a-9c8b-1f2d3c4b5a6e",
  "email": "seller@example.com",
  "role": "SELLER"
}
\end{lstlisting}

\hrulefill

\section{BundleController}
\textbf{Base path:} \texttt{/api/bundles}

\subsection{List available bundles (paged)}
\texttt{GET /api/bundles}

\paragraph{Example Response:}
\begin{itemize}
    \item Returns \texttt{Page<BundlePosting>}
\end{itemize}

\subsection{Get bundle by id}
\texttt{GET /api/bundles/\{id\}}

\paragraph{Example Response:}
\begin{itemize}
    \item \texttt{BundlePosting}
\end{itemize}

\subsection{Get bundles by seller id}
\texttt{GET /api/bundles/seller/\{sellerId\}}

\paragraph{Example Response:}
\begin{itemize}
    \item \texttt{List<BundlePosting>}
\end{itemize}

\subsection{Create bundle posting}
\texttt{POST /api/bundles}

\paragraph{Example Request:}
\begin{lstlisting}[language=json]
{
  "categoryId": "2e5df1b0-1e2c-4a1b-8d71-3d9c5f5d4f90",
  "pickupStartAt": "2026-01-30T10:00:00Z",
  "pickupEndAt": "2026-01-30T18:00:00Z",
  "quantityTotal": 10,
  "quantityReserved": 1,
  "priceCents": 899,
  "discountPct": 30,
  "contentsText": "Assorted bakery items",
  "allergensText": "Gluten",
  "estimatedWeightGrams": 1200,
  "activate": true
}
\end{lstlisting}

\paragraph{Example Response:}
\begin{itemize}
    \item Saved \texttt{BundlePosting}
\end{itemize}

\paragraph{Notes:}
\begin{itemize}
    \item Creates \texttt{BundlePosting} with \texttt{quantityReserved = 0}
\end{itemize}

\subsection{Update bundle (seller)}
\texttt{PUT /api/bundles/\{id\}}

\paragraph{Example Request:}
\begin{lstlisting}[language=json]
{
  "quantityTotal": 12,
  "priceCents": 799,
  "discountPct": 40,
  "contentsText": "Bakery items and sandwiches",
  "allergensText": "Nuts"
}
\end{lstlisting}

\paragraph{Example Response:}
\begin{itemize}
    \item Updated \texttt{BundlePosting}
\end{itemize}

\subsection{Activate bundle}
\texttt{POST /api/bundles/\{id\}/activate}

\paragraph{Example Response:}
\begin{itemize}
    \item Updated \texttt{BundlePosting}
\end{itemize}

\subsection{Close bundle}
\texttt{POST /api/bundles/\{id\}/close}

\paragraph{Example Response:}
\begin{itemize}
    \item Updated \texttt{BundlePosting}
\end{itemize}

\hrulefill

\section{CategoryController}
\textbf{Base path:} \texttt{/api/categories}

\subsection{List all categories}
\texttt{GET /api/categories}

\paragraph{Example Response:}
\begin{itemize}
    \item \texttt{List<Category>}
\end{itemize}

\paragraph{Example Response (JSON):}
\begin{lstlisting}[language=json]
[
  { "categoryId": "11111111-1111-1111-1111-111111111111", "name": "Bakery" },
  { "categoryId": "22222222-2222-2222-2222-222222222222", "name": "Hot Meals" }
]
\end{lstlisting}

\hrulefill

\section{GamificationController}
\textbf{Base path:} \texttt{/api/gamification}

\subsection{Get an employee's streak stats}
\texttt{GET /api/gamification/streak/\{employeeId\}}

\paragraph{Example Response:}
\begin{lstlisting}[language=json]
{
  "currentStreakWeeks": 3,
  "bestStreakWeeks": 7,
  "lastRescueWeekStart": "2026-01-26"
}
\end{lstlisting}

\subsection{Get an employee's impact summary}
\texttt{GET /api/gamification/impact/\{employeeId\}}

\paragraph{Example Response:}
\begin{lstlisting}[language=json]
{
  "totalRescues": 12,
  "totalMealsSaved": 12,
  "totalCo2eSavedKg": 4.8,
  "currentStreakWeeks": 3,
  "badgesCount": 2
}
\end{lstlisting}

\paragraph{Notes:}
\begin{itemize}
    \item \texttt{totalRescues} values from \texttt{rescueEventRepo.countByEmployee\_EmployeeId}
    \item \texttt{totalMealsSaved} values from \texttt{rescueEventRepo.sumMealsByEmployee}
    \item \texttt{totalCo2eSavedKg} values from \texttt{rescueEventRepo.sumCo2eByEmployee / 1000.0}
\end{itemize}

\subsection{List employee's badges}
\texttt{GET /api/gamification/badges/\{employeeId\}}

\paragraph{Example Response:}
\begin{itemize}
    \item \texttt{List<EmployeeBadge>}
\end{itemize}

\subsection{List all badges}
\texttt{GET /api/gamification/badges}

\paragraph{Example Response:}
\begin{itemize}
    \item \texttt{List<Badge>}
\end{itemize}

\hrulefill

\section{IssueController}
\textbf{Base path:} \texttt{/api/issues}

\subsection{List all issues for a seller}
\texttt{GET /api/issues/seller/\{sellerId\}}

\paragraph{Example Response:}
\begin{itemize}
    \item \texttt{List<IssueReport>}
\end{itemize}

\subsection{List open issues for a seller}
\texttt{GET /api/issues/seller/\{sellerId\}/open}

\paragraph{Example Response:}
\begin{itemize}
    \item \texttt{List<IssueReport>}
\end{itemize}

\subsection{Create issue}
\texttt{POST /api/issues}

\paragraph{Example Request:}
\begin{lstlisting}[language=json]
{
  "postingId": "aaaaaaaa-aaaa-aaaa-aaaa-aaaaaaaaaaaa",
  "reservationId": "bbbbbbbb-bbbb-bbbb-bbbb-bbbbbbbbbbbb",
  "employeeId": "cccccccc-cccc-cccc-cccc-cccccccccccc",
  "type": "QUALITY",
  "description": "Item quality was not as described."
}
\end{lstlisting}

\paragraph{Example Response:}
\begin{itemize}
    \item Saved \texttt{IssueReport}
\end{itemize}

\subsection{Seller respond}
\texttt{POST /api/issues/\{id\}/respond}

\paragraph{Example Request:}
\begin{lstlisting}[language=json]
{
  "resolve": true
}
\end{lstlisting}

\paragraph{Example Response:}
\begin{itemize}
    \item Updated \texttt{IssueReport}
\end{itemize}

\paragraph{Notes:}
\begin{itemize}
    \item Updates status to \texttt{RESPONDED}.
    \item If \texttt{resolve = true}, sets status to \texttt{RESOLVED} and timestamps \texttt{resolvedAt}.
\end{itemize}

\subsection{Resolve an issue}
\texttt{POST /api/issues/\{id\}/resolve}

\paragraph{Example Response:}
\begin{itemize}
    \item Updated \texttt{IssueReport}
\end{itemize}

\paragraph{Notes:}
\begin{itemize}
    \item Sets status to \texttt{RESOLVED}.
    \item Sets \texttt{resolvedAt = Instant.now()}.
\end{itemize}

\hrulefill

\section{ReservationController}
\textbf{Base path:} \texttt{/api/reservations}

\subsection{List reservations by organisation}
\texttt{GET /api/reservations/org/\{orgId\}}

\paragraph{Example Response:}
\begin{itemize}
    \item \texttt{List<Reservation>}
\end{itemize}

\subsection{List reservations by employee}
\texttt{GET /api/reservations/employee/\{employeeId\}}

\paragraph{Example Response:}
\begin{itemize}
    \item \texttt{List<Reservation>}
\end{itemize}

\subsection{Create a reservation (generates claim code)}
\texttt{POST /api/reservations}

\paragraph{Example Request:}
\begin{lstlisting}[language=json]
{
  "postingId": "aaaaaaaa-aaaa-aaaa-aaaa-aaaaaaaaaaaa",
  "organisationId": "dddddddd-dddd-dddd-dddd-dddddddddddd",
  "employeeId": "cccccccc-cccc-cccc-cccc-cccccccccccc",
  "quantity": 1
}
\end{lstlisting}

\paragraph{Example Response:}
\begin{lstlisting}[language=json]
{
  "reservationId": "bbbbbbbb-bbbb-bbbb-bbbb-bbbbbbbbbbbb",
  "claimCode": "A1B2C3D4E5",
  "claimCodeLast4": "D4E5",
  "pickupStartAt": "2026-01-28T10:00:00Z",
  "pickupEndAt": "2026-01-28T12:00:00Z",
  "sellerName": "ABC",
  "locationText": "Exeter"
}
\end{lstlisting}

\paragraph{Notes:}
\textbf{Claim code generation process:}
\begin{itemize}
    \item Returns plaintext \texttt{claimCode} in response.
    \item Stores hash as \texttt{claimCodeHash}.
    \item Stores \texttt{claimCodeLast4}.
\end{itemize}

\subsection{Verify claim code}
\texttt{POST /api/reservations/\{id\}/verify}

\paragraph{Example Request:}
\begin{lstlisting}[language=json]
{
  "claimCode": "A1B2C3D4E5"
}
\end{lstlisting}

\paragraph{Example Response:}
\begin{lstlisting}[language=json]
{
  "success": true,
  "message": "Bundle collected successfully"
}
\end{lstlisting}

\subsection{Mark reservation as no-show}
\texttt{POST /api/reservations/\{id\}/no-show}

\paragraph{Example Response:}
\begin{itemize}
    \item Updated \texttt{Reservation}
\end{itemize}

\paragraph{Notes:}
\begin{itemize}
    \item Sets status to \texttt{NO\_SHOW} and sets \texttt{noShowMarkedAt}.
    \item Sets \texttt{bundle.quantityReserved -= 1}.
\end{itemize}

\subsection{Cancel reservation (only if RESERVED)}
\texttt{POST /api/reservations/\{id\}/cancel}

\paragraph{Example Response:}
\begin{itemize}
    \item Updated \texttt{Reservation}
\end{itemize}

\paragraph{Notes:}
\begin{itemize}
    \item Sets status to \texttt{CANCELLED}.
    \item Sets \texttt{bundle.quantityReserved -= 1}.
\end{itemize}

\subsection{Assign an employee to a reservation}
\texttt{POST /api/reservations/\{id\}/assign/\{employeeId\}}

\paragraph{Example Response:}
\begin{itemize}
    \item Updated \texttt{Reservation}
\end{itemize}

\paragraph{Notes:}
\begin{itemize}
    \item Checks reservation and employee exists.
    \item Sets \texttt{reservation.employee = employee}.
\end{itemize}

\end{document}


